\section{Conclusions}
\label{sec:Conclusions}
The IoT-Sim is a practical, lightweight, and flexible platform for designing, configuring, and exploring IoT network topologies while testing attack-detection models under controlled and reproducible conditions. Its modular client-side architecture supports rapid prototyping and experimentation without the need for a specialized infrastructure, allowing researchers and students to generate realistic traffic, inject diverse attack behaviors, and observe model responses in real time.

The integrated support for importing external scripts and TensorFlow models facilitates a direct comparison of heterogeneous detection approaches within identical simulated environments. This capability promotes rigorous benchmarking of algorithms, hyperparameter settings, and deployment strategies and strengthens the reproducibility of experimental results by allowing the same scenarios and artifacts to be shared and re-executed.

The simulator focuses on usability and configurability, reducing the entry barrier for security experimentation in resource-constrained IoT environments, and promotes iterative exploration of model robustness across variable network conditions, such as latency and bandwidth limitations. The visual component-based interface and the per-connection model attachment mechanism facilitate the precise analysis of detection performance and enable investigations into localized versus global detection strategies.

The tool’s current implementation and demonstrated use in laboratory comparisons expose research opportunities for lightweight distributed detection, model adaptation to changing traffic patterns, and privacy-preserving training workflows for IoT ecosystems. The simulator creates a controlled testbed for studying resilience to evasion techniques, sensitivity to attack parameterization, and trade-offs between detection accuracy, latency, and computational overhead on constrained devices.

Future development directions will enhance realism and extend applicability by incorporating additional communication protocols, richer traffic generators, and support for longitudinal and large-scale scenario orchestration. These enhancements will increase the fidelity of experiments, broaden the range of empirical questions that can be addressed, and strengthen the simulator’s role as a reference platform for the design and evaluation of intelligent, deployable IoT security solutions.

Overall, the IoT-Sim constitutes a valuable contribution to the research infrastructure for IoT cybersecurity by combining accessibility, modularity, and machine learning integration to accelerate reproducible experimentation, methodological comparison, and the development of more robust intrusion detection strategies for heterogeneous IoT environments.
